% ============ 第五章 密度问题 (原书 pp.188-205) ============
\chapter{密度问题}\label{ch:5}

我们在上一章介绍了复变函数的解析性质(Euler 乘积展开、解析延拓、函数方程、零点和极点特性等). 在历史上, 这些函数的引进以及对解析性质的研究都具有数论目的. Riemann zeta 函数 $\zeta(s)$ 用来研究素数的分布, Dirichlet $L$-函数 $L(s,\chi)$ 用来研究算术级数中素数的分布, 而 Dedekind zeta 函数 $\zeta_K(s)$ 则用来研究数域 $K$ 中具有各种性质的素理想的分布. 

\begin{sloppypar}
公元前 $3$ 世纪欧几里德证明了素数有无穷多个, 证明是很初等的, 只用到算术基本定理: 如果只有有限多个素数 $p_1,p_2,\cdots,p_n$, 考虑正整数 $N=p_1p_2\cdots p_n+1$, 它应当是一些素数的乘积, 但是素数只有 $p_1,p_2,\cdots,p_n$, 而它们均不是 $N$ 的因子, 这就导致矛盾. 如果我们定义数论函数
\[
  f(n)=\begin{cases}
    1, & \text{若 } n\ \text{为素数},\\
    0, & \text{否则},
  \end{cases}
\]
则函数
\[
  \pi(x)=\sum_{n\leqslant x}f(n)=\sum_{p\leqslant x}1\qquad(x\in\R)
\]
恰好是不超过 $x$ 的素数个数. 素数有无穷多个相当于 $\pi(x)\to+\infty$(当 $x\to+\infty$ 时). 进而我们要问: 素数在所有正整数中间的分布如何？占有多大比例？首先, 对每个正整数 $n\geqslant 2$, 数 $(n!)+2$, $(n!)+3,\cdots,(n!)+n$ 都不是素数, 所以我们总有两个相距很远的素数, 它们中间不再有素数. 另一方面, 人们猜想有无限多个素数 $p$, 使得 $p+2$ 也是素数, 即两个素数距离只有 $2$(这个猜想至今未解决). 这表明素数在正整数中间的分布是很不规则的. 不超过 $x$ 的正整数共有 $[x]$ 个, $|x-[x]|\leqslant 1$. 而不超过 $x$ 的素数共有 $\pi(x)$ 个. 于是 $\frac{\pi(x)}{[x]}$ 就是区间 $(1,x)$ 之内素数在正整数中所占的比例. 利用古典的希腊筛法可证明
\[
  \lim_{x\to\infty}\frac{\pi(x)}{[x]}=\lim_{x\to\infty}\frac{\pi(x)}{x}=0.
\]
这意味着素数在正整数中是稀疏的. $18$ 世纪至 $19$ 世纪, Legendre 和 Gauss 等人基于大量手算结果猜想
\begin{equation}  \lim_{x\to\infty}\frac{\pi(x)}{\frac{x}{\log x}}=1\tag{1}\end{equation}
(今后 $\log x$ 均指是自然对数). 这个猜想由两个法国数学家 J. Hadamard 和 C. J. de la Vallée Poussin 于 $1896$ 年各自独立地证明, 方法则完全是解析的, 即利用 $\zeta(s)$ 和相当深刻的复变函数理论. 直到 $1949$ 年, 又有两个数学家 A. Selberg 和 P. Erd\H{o}s 独立地给出不用复变函数理论的初等证明. 
\end{sloppypar}

接下来人们考虑算术级数(即等差数列)中是否包含无穷多个素数. 设 $k$ 和 $l$ 为正整数, $k\geqslant 2$, 而 $1\leqslant l<k$. 则算术级数 $l,l+k,l+2k,\cdots$ 中每个数都有因子 $(l,k)$. 所以当 $(l,k)\geqslant 2$ 时, 这个算术级数当中至多有一个素数. 所以我们要假设 $(l,k)=1$. 对于固定的公差 $k\geqslant 2$, 满足 $1\leqslant l<k$, $(l,k)=1$ 的 $l$ 共有 $\varphi(k)$ 个. 我们要问: 这 $\varphi(k)$ 个算术级数 $\{l,l+k,l+2k,\cdots\}=\{n\geqslant 1:n\equiv l\ (\mathrm{mod}\ k)\}$ 中是否都有无限多个素数？对于一些特殊情形我们有初等解决方法. 例如对 $k=4$ 和 $l=3$ 的情形, 可以用欧几里德方法证明满足 $p\equiv 3\ (\mathrm{mod}\ 4)$ 的素数 $p$ 有无穷多个. 对一般情形, Dirichlet 引进了 $L$-函数 $L(s,\chi)$, 其中 $\chi$ 是模 $k$ 的特征. 利用 $L(s,\chi)$ 的解析性质可以肯定地回答上述问题. 确切地说, 对于 $k\geqslant 2$, $(l,k)=1$, $1\leqslant l<k$, 我们用 $\pi(x;k,l)$ 表示满足 $p\leqslant x$, $p\equiv l\ (\mathrm{mod}\ k)$ 的素数 $p$ 的个数. 可以证明: 
\begin{equation}  \lim_{x\to\infty}\frac{\pi(x;k,l)}{\frac{x}{\log x}}=\frac{1}{\varphi(k)}.\tag{2}\end{equation}
由此特别得到 $\pi(x;k,l)\to+\infty$(当 $x\to\infty$ 时), 即算术级数 $A(k,l)=\{n\geqslant 1:n\equiv l\ (\mathrm{mod}\ k)\}$ 中有无穷多素数. 并且 (2) 式表明对固定的 $k\geqslant 2$, 素数是平均地分配在 $\varphi(k)$ 个算术级数 $A(k,l)$($(l,k)=1$, $1\leqslant l<k$)之中. $1920$ 年, 德国数学家 Landau 利用 Dedekind zeta 函数 $\zeta_K(s)$ 和同样的解析方法, 证明了: 对每个数域 $K$, 以 $\pi_K(x)$ 表示 $O_K$ 中范不超过 $x$ 的素理想个数, 则
\begin{equation}  \lim_{x\to\infty}\frac{\pi_K(x)}{\frac{x}{\log x}}=1.\tag{3}\end{equation}

\begin{sloppypar}
我们还需要研究具有特殊性质的素理想的分布问题. 比如设 $L/K$ 为数域的伽罗瓦扩张. 由判别式定理知道 $O_K$ 中素理想 $\mathfrak{p}$ 在 $L$ 中分歧的只有有限多个. 如果 $\mathfrak{p}$ 在 $L$ 中不分歧, 则 $\mathfrak{p}O_L=\mathfrak{P}_1\mathfrak{P}_2\cdots\mathfrak{P}_g$, $N_{L/K}(\mathfrak{P}_i)=\mathfrak{p}^f$($1\leqslant i\leqslant g$), $gf=[L:K]$. 我们的问题是: 对固定的 $(g,f)$, 上述分解模式的 $\mathfrak{p}$ 是否有无穷多个？这类问题的解决均采用解析方法. 本章中介绍如何用解析方法来解决这类问题, 以及这些结果的数论应用. 
\end{sloppypar}

% ============ 第五章 §1 Dirichlet 密度 (原书 pp.190-196) ============
\section{Dirichlet 密度}\label{sec:ch5-s1}

我们先从一个简单的例子讲起, 看一下如何利用 Dirichlet $L$-函数 $L(s,\chi)$ 的解析特性来证明算术级数中有无穷多个素数. 以下设
\[
  k\geqslant 2,\quad 1\leqslant l<k,\quad (l,k)=1,
\]
$A(k,l)$ 表示算术级数 $\{l+kt\mid t\geqslant 0\}=\{n>0:n\equiv l\ (\mathrm{mod}\ k)\}$, $\pi(x;k,l)=\sum\limits_{\substack{p\leqslant x\\ p\equiv l\ (\mathrm{mod}\ k)}}1$ 表示集合 $A(k,l)$ 中不超过 $x$ 的素数个数. 

\begin{theorem}[Dirichlet]\label{thm:5.1}
设 $k\geqslant 2$, $1\leqslant l<k$, $(l,k)=1$, 则算术级数 $A(k,l)$ 中有无穷多个素数. 
\end{theorem}

\begin{proof}
以 $X$ 表示所有模 $k$ 的 Dirichlet 特征组成的集合. 对每个 $\chi\in X$, 当 $\operatorname{Re}(s)>1$ 时, 
\[
  \log L(s,\chi)=\sum\limits_p\log(1-\chi(p)p^{-s})^{-1}
  =\sum\limits_p\sum\limits_{m=1}^{\infty}\frac{\chi(p)^m}{m}p^{-ms}
  =\sum\limits_p\chi(p)p^{-s}+g_\chi(s),
\]
其中
\[
  g_\chi(s)=\sum\limits_p\sum\limits_{m=2}^{\infty}\frac{\chi(p)^m}{m}p^{-ms}.
\]
易知级数 $g_\chi(s)$ 在 $\operatorname{Re}(s)>\frac{1}{2}$ 中收敛, 从而 $g_\chi(s)$ 在 $s=1$ 处连续. 另一方面, 由特征的正交性质可知
\begin{align*}
  \sum\limits_{p\equiv l\ (\mathrm{mod}\ k)}p^{-s}
  &=\frac{1}{\varphi(k)}\sum\limits_{\chi\in X}\overline{\chi}(l)\sum\limits_p\chi(p)p^{-s}\\
  &=\frac{1}{\varphi(k)}\left(\sum\limits_{\chi\in X}\overline{\chi}(l)\log L(s,\chi)-\sum\limits_\chi\overline{\chi}(l)g_\chi(s)\right).
\end{align*}
如果 $\chi\neq\chi_0$, 则 $L(s,\chi)$ 在 $s=1$ 处正则. 而
\[
  \log L(s,\chi_0)\sim\log\zeta(s)\sim-\log(s-1)\qquad(\text{当 } s\to1^+\ \text{时}).
\]
此外, $\sum\limits_\chi\overline{\chi}(l)g_\chi(s)$ 在 $s=1$ 处正则. 因此
\[
  \sum\limits_{p\in A(k,l)}p^{-s}=\sum\limits_{p\equiv l\ (\mathrm{mod}\ k)}p^{-s}
  \sim\frac{1}{\varphi(k)}(-\log(s-1))\qquad(\text{当 } s\to1^+\ \text{时}).
\]
这就表明当 $s\to1^+$ 时 $\sum\limits_{p\in A(k,l)}p^{-s}\to+\infty$. 特别地, $A(k,l)$ 中包含无限多个素数. 证毕. 
\end{proof}

这个定理的证明给出
\begin{equation}  \lim\limits_{s\to1^+}\frac{\sum\limits_{p\in A(k,l)}p^{-s}}{-\log(s-1)}=\frac{1}{\varphi(k)},\tag{4}\end{equation}
由此启发出如下的一个概念: 

\begin{definition}\label{def:ch5-1}
设 $K$ 是代数数域, $A$ 是 $O_K$ 的一个素理想集合. 如果极限
\[
  \lim\limits_{s\to1^+}\sum\limits_{\mathfrak{p}\in A}N(\mathfrak{p})^{-s}\big/-\log(s-1)
\]
存在, 称此极限为素理想集合 $A$ 的 \emph{Dirichlet 密度}（简称作 $D$-密度）, 表示成 $\delta(A)$. 
\end{definition}

\begin{example}
取 $K=\Q$, 则整数环 $\Z$ 中的非零素理想和素数是一一对应的. 由 (4) 式可知对每个 $k\geqslant 2$, $1\leqslant l<k$, $(l,k)=1$, 满足 $p\equiv l\ (\mathrm{mod}\ k)$ 的素数 $p$ 组成的集合, 其 $D$-密度为 $\frac{1}{\varphi(k)}$. 
\end{example}

由 $D$-密度的定义立刻推出如下性质: 

\begin{lemma}\label{lem:ch5-1}
设 $K$ 是代数数域, $A$ 和 $A'$ 是 $O_K$ 的两个素理想集合. 
\begin{enumerate}[(1)]
  \item 如果 $A$ 是有限集合, 则 $\delta(A)=0$. 换句话说, 如果 $\delta(A)$ 存在并且 $\delta(A)>0$, 则 $A$ 为无限集合. 
  \item 若 $A\subseteq A'$, 并且 $\delta(A)$ 和 $\delta(A')$ 均存在, 则 $\delta(A)\leqslant\delta(A')$. 
  \item 若 $\delta(A)$ 存在, 则 $0\leqslant\delta(A)\leqslant 1$. 
  \item 若 $A\cap A'=\varnothing$, 并且 $\delta(A)$ 和 $\delta(A')$ 均存在, 则 $\delta(A\cup A')$ 也存在, 并且 $\delta(A\cup A')=\delta(A)+\delta(A')$. 
  \item 记 $A\Delta A'=(A\cup A')-(A\cap A')$. 如果 $\delta(A)$ 存在并且 $\delta(A\Delta A')=0$, 则 $\delta(A')$ 和 $\delta(A\cap A')$ 也存在并且均等于 $\delta(A)$. 
\end{enumerate}
\end{lemma}

\begin{proof}
(1) 若 $A$ 是有限集合, 则 $\sum\limits_{\mathfrak{p}\in A}N\mathfrak{p}^{-s}$ 有限. 但是当 $s\to1^+$ 时 $-\log(s-1)\to+\infty$. 所以 $\delta(A)=\lim\limits_{s\to1^+}\sum\limits_{\mathfrak{p}\in A}N\mathfrak{p}^{-s}/-\log(s-1)=0$. 

(2) 由于 $\sum\limits_{\mathfrak{p}\in A}N\mathfrak{p}^{-s}=\sum\limits_{\mathfrak{p}\in A'}N\mathfrak{p}^{-s}-\sum\limits_{\mathfrak{p}\in A'-A}N\mathfrak{p}^{-s}\leqslant\sum\limits_{\mathfrak{p}\in A'}N\mathfrak{p}^{-s}$（当 $s>1$ 时）. 从而 $\delta(A)\leqslant\delta(A')$. 

(3) 显然 $\delta(A)\geqslant 0$. 另一方面, 以 $P$ 表示 $O_K$ 中全体素理想组成的集合, 则由 (2) 可知 $\delta(A)\leqslant\delta(P)$. 我们现在证明 $\delta(P)=1$. 我们在上节已经证明 Dedekind zeta 函数
\[
  \zeta_K(s)=\prod_{\mathfrak{p}\in P}(1-N\mathfrak{p}^{-s})^{-1}
\]
在 $s=1$ 处有单极点, 并且 $\zeta_K(s)\sim\rho_K(s-1)^{-1}$（$\rho_K\neq 0$）（当 $s\to1^+$ 时）. 从而 $\log\zeta_K(s)\sim-\log(s-1)$. 但是
\begin{align*}
  \log\zeta_K(s)&=\sum\limits_{\mathfrak{p}\in P}\log(1-N\mathfrak{p}^{-s})^{-1}
  =\sum\limits_{\mathfrak{p}\in P}\sum\limits_{m=1}^{\infty}\frac{N\mathfrak{p}^{-ms}}{m}\\
  &=\sum\limits_{\mathfrak{p}\in P}N\mathfrak{p}^{-s}+\sum\limits_{\mathfrak{p}\in P}\sum\limits_{m=2}^{\infty}\frac{N\mathfrak{p}^{-ms}}{m},
\end{align*}
而右边第 $2$ 项当 $s\to1^+$ 时收敛, 因此, 
\[
  \lim\limits_{s\to1^+}\frac{\sum\limits_{\mathfrak{p}\in P}N\mathfrak{p}^{-s}}{-\log(s-1)}
  =\lim\limits_{s\to1^+}\frac{\log\zeta_K(s)}{-\log(s-1)}=1.
\]
即 $\delta(P)=1$. 

(4) 当 $A\cap A'=\varnothing$ 时, $\sum\limits_{\mathfrak{p}\in A}N\mathfrak{p}^{-s}+\sum\limits_{\mathfrak{p}\in A'}N\mathfrak{p}^{-s}=\sum\limits_{\mathfrak{p}\in A\cup A'}N\mathfrak{p}^{-s}$, 由此即知 $\delta(A\cup A')=\delta(A)+\delta(A')$. 

(5) 由于 $A\cap A'\subseteq A$ 并且 $A-(A\cap A')\subseteq A\Delta A'$. 根据 $D$-密度的定义和 $\delta(A\Delta A')=0$, 易知 $\delta(A\cap A')$ 存在并且等于 $\delta(A)$. 再由 $A\cap A'\subseteq A'$ 并且 $A'-(A\cap A')\subseteq A\Delta A'$, 又知 $\delta(A')$ 存在并且等于 $\delta(A\cap A')=\delta(A)$. 
\end{proof}

\begin{remark}
由引理 \ref{lem:ch5-1} 的证明可知, $O_K$ 中所有素理想组成的集合其 $D$-密度为 $1$. 注意前面的 (3) 式是更强的结果. 因为若 $A$ 是 $O_K$ 的一个素理想集合, 令 $\pi(x,A)$ 表示 $A$ 中范不超过 $x$ 的素理想的个数. 如果 $\lim\limits_{x\to\infty}\frac{\pi(x,A)}{x/\log x}=C(A)$, 则可以证明 $A$ 存在 $D$-密度并且 $\delta(A)=C(A)$. （证明从略）. 要想证明前面的 (1)、(2)、(3) 式, 还需要更精细的分析工具. 
\end{remark}

现在再给出计算 $D$-密度的新的例子. 

\begin{theorem}\label{thm:5.2}
设 $L/K$ 是数域的伽罗瓦扩张, $n=[L:K]$. 以 $A$ 表示 $O_K$ 中在 $L$ 完全分裂的那些素理想 $\mathfrak{p}$ 组成的集合, 则 $\delta(A)=\frac{1}{n}$. 
\end{theorem}

\begin{proof}
当 $s>1$ 时
\begin{equation}\tag{5}
  \log\zeta_L(s)=\sum\limits_{\mathfrak{P}}N_{L/\Q}(\mathfrak{P})^{-s}+O(1)
  =\sum\limits_{\mathfrak{p}}\sum\limits_{\mathfrak{P}\mid\mathfrak{p}}N_{L/\Q}(\mathfrak{P})^{-s}+O(1).
\end{equation}
由于只有有限个 $\mathfrak{p}$ 在 $L$ 中分歧, 它们不影响对 $\delta(A)$ 的计算. 从而将它们排除之后, 我们有 $\mathfrak{p}=\mathfrak{P}_1\cdots\mathfrak{P}_g$, $gf=n$. 当 $f\geqslant 2$ 时, $N_{L/\Q}(\mathfrak{P})^{-s}=N_{K/\Q}(\mathfrak{p})^{-fs}\leqslant N_{K/\Q}(\mathfrak{p})^{-2s}$. (5) 式右边对应于 $f\geqslant 2$ 的那些项在 $s\geqslant 1$ 时一致收敛. 而当 $f=1$（即 $\mathfrak{p}\in A$）时, $\mathfrak{p}$ 有 $n$ 个因子 $\mathfrak{P}$, $N_{L/\Q}(\mathfrak{P})=N_{K/\Q}(\mathfrak{p})$. 因此, 
\[
  \log\zeta_L(s)=\sum\limits_{\mathfrak{p}\in A}\sum\limits_{\mathfrak{P}\mid\mathfrak{p}}N_{L/\Q}(\mathfrak{P})^{-s}+O(1)
  =n\sum\limits_{\mathfrak{p}\in A}N_{K/\Q}(\mathfrak{p})^{-s}+O(1).
\]
另一方面, $\log\zeta_L(s)=-\log(s-1)+O(1)$（$s\to1^+$）. 于是
\[
  \lim\limits_{s\to1^+}\frac{\sum\limits_{\mathfrak{p}\in A}N_{K/\Q}(\mathfrak{p})^{-s}}{-\log(s-1)}
  =\lim\limits_{s\to1^+}\frac{\frac{1}{n}\log\zeta_L(s)+O(1)}{\log\zeta_L(s)+O(1)}=\frac{1}{n}.
\]
即 $\delta(A)=\frac{1}{n}$. 
\end{proof}

如果 $L/K$ 不是伽罗瓦扩张, 我们记 $N$ 为 $L$ 在 $K$ 上的伽罗瓦闭包, 即包含 $L$ 的 $K$ 之最小伽罗瓦扩域. 如果 $L=K(\alpha)$, 并且 $f(x)\in K[x]$ 是 $\alpha$ 在 $K$ 上的极小多项式, 则 $N$ 就是将 $f(x)$ 的所有根添加到 $K$ 上得到的域. 也就是说, $N$ 是 $L$ 的所有 $K$-共轭域的并. 

\begin{sloppypar}
\begin{corollary}[1]
设 $L/K$ 是数域的扩张, $N$ 是 $L$ 在 $K$ 上的正规闭包, 令 $A=\{O_K\ \text{中素理想}\ \mathfrak{p}\mid\mathfrak{p}\ \text{在}\ L\ \text{中完全分裂}\}$, 则
\[
  \delta(A)=1/[N:K].
\]
\end{corollary}
\end{sloppypar}

\begin{proof}
令 $A'=\{O_K\ \text{中素理想}\ \mathfrak{p}\mid\mathfrak{p}\ \text{在}\ N\ \text{中完全分裂}\}$. 由于 $K\subseteq L\subseteq N$, 可知对每个 $\sigma\in G=\Gal(N/K)$ 均有 $K\subseteq\sigma(L)\subseteq N$. 现设 $\mathfrak{p}\in A$, 令 $\mathfrak{P}$ 为 $N$ 中素理想, $\mathfrak{P}\mid\mathfrak{p}$. 由于 $\mathfrak{p}$ 在 $L$ 中完全分裂, 从而关于 $\mathfrak{P}$ 的分解域 $K_D\supseteq L$（为什么？）. 同样地, 对于每个 $\sigma\in G$, $\mathfrak{p}$ 在 $\sigma(L)$ 中也是完全分裂的（为什么？）. 因此 $K_D\supseteq\sigma(L)$. 于是 $N\supseteq K_D\supseteq\bigcup_{\sigma\in G}\sigma(L)=N$, 从而 $K_D=N$. 这表明 $\mathfrak{p}$ 在 $N$ 中完全分裂, 即 $\mathfrak{p}\in A'$. 反之, 若 $\mathfrak{p}\in A'$, 即 $\mathfrak{p}$ 在 $N$ 中完全分裂, 从而在 $L$ 中显然也完全分裂, 因此 $\mathfrak{p}\in A$, 这就表明 $A=A'$. 然后由定理 \ref{thm:5.2} 即知 $\delta(A)=\delta(A')=1/[N:K]$. 
\end{proof}

设 $K$ 为数域, $A$ 和 $A'$ 是 $O_K$ 的两个素理想集合. 如果 $\delta(A\Delta A')=0$, 即 $A$ 和 $A'$ 彼此不同的素理想组成 $D$-密度为 $0$ 的集合, 则称 $A$ 和 $A'$ \emph{几乎相等}. 引理 \ref{lem:ch5-1} 表明, 若 $A$ 和 $A'$ 几乎相等, 如果 $A$, $A'$, $A\cap A'$ 和 $A\cup A'$ 当中任何一个存在 $D$-密度, 则其余 $3$ 个也存在 $D$-密度, 并且这 $4$ 个集合的 $D$-密度相等（这是由于若 $A$ 和 $A'$ 几乎相等, 则 $A$, $A'$, $A\cap A'$ 和 $A\cup A'$ 当中任意两个都几乎相等）. 

\begin{corollary}[2]
设 $L_1/K$ 和 $L_2/K$ 均是数域的伽罗瓦扩张. 以 $S_i$ 表示在 $L_i$ 中完全分裂的 $O_K$ 中素理想构成的集合（$i=1,2$）. 如果 $S_1$ 和 $S_2$ 几乎相等, 则 $L_1=L_2$. 
\end{corollary}

\begin{proof}
令 $L=L_1L_2$, 则 $L/K$ 也是伽罗瓦扩张. 以 $S$ 表示在 $L$ 完全分裂的 $O_K$ 中素理想构成的集合. 我们在第二章证明了: $O_K$ 的素理想 $\mathfrak{p}$ 在 $L$ 中完全分裂当且仅当 $\mathfrak{p}$ 在 $L_1$ 和 $L_2$ 中均完全分裂. 这意味着 $S=S_1\cap S_2$. 由于假设 $S_1$ 和 $S_2$ 几乎相等, 并且定理 \ref{thm:5.2} 表明 $\delta(S_i)=\frac{1}{[L_i:K]}$（$i=1,2$）, 可知 $\delta(S)=\delta(S_i)$（$i=1,2$）即 $[L:K]=\delta(S)^{-1}=\delta(S_i)^{-1}=[L_i:K]$（$i=1,2$）. 但是 $L\supseteq L_i$, 于是 $L=L_i$（$i=1,2$）. 即 $L_1=L=L_2$. 证毕. 
\end{proof}

\begin{remark}
定理 \ref{thm:5.2} 表明, 如果 $L/K$ 是数域的伽罗瓦扩张, 则 $K$ 中存在相当多（$D$-密度为 $1/[L:K]$）的素理想 $\mathfrak{p}$ 在 $L$ 中完全分裂. 而推论 2 表明, 这些完全分裂的素理想组成的集合也完全决定了伽罗瓦扩域 $L$. 
\end{remark}

\begin{sloppypar}
现在我们将定理 \ref{thm:5.2} 及其推论用于有限域上多项式的因子分解问题. 设 $L/K$ 是数域的扩张, $L=K(\theta)$, $\theta\in O_L$, $f(x)\in O_K[x]$ 是 $\theta$ 在 $K$ 上的极小多项式, 则 $f(x)$ 是不可约的首 $1$ 多项式, 并且 $\deg f(x)=n=[L:K]$. 我们在第二章中看到, 除了有限多个可能之外, 对于 $K$ 中每个在 $L$ 中不分歧的素理想 $\mathfrak{p}$, $\mathfrak{p}O_L$ 在 $O_L$ 中的分解型式（即 $\mathfrak{p}O_L=\mathfrak{P}_1\cdots\mathfrak{P}_g$ 的参数 $g$ 和 $f_i=f(\mathfrak{P}_i\mid\mathfrak{p})$（$1\leqslant i\leqslant g$）, $\sum\limits_{i=1}^{g}f_i=n$）与多项式 $f(x)$ 在有限域 $O_K/\mathfrak{p}$ 上的分解型式是一样的. 于是由定理 \ref{thm:5.2} 中关于 $\mathfrak{p}$ 的完全分裂性结果可以推得 $f(x)\ (\mathrm{mod}\ \mathfrak{p})$ 的完全分裂结果. 为简单起见我们只考虑 $K=\Q$ 的情形. 虽然对于一般情形其结果也是对的. 
\end{sloppypar}

\setcounter{definition}{2}% 原书定义编号跳号（无定义2）

\begin{definition}\label{def:ch5-3}
设 $f(x)\in\Z[x]$ 是 $\Z$ 上的首 $1$ 不可约 $n$ 次多项式, $p$ 为素数. 如果 $f(x)$ 在 $p$ 元域 $\Z/p\Z$ 上分解成 $n$ 个不同的一次因子之积（这相当于说 $f(x)$ 在 $\Z/p\Z$ 中有 $n$ 个不同的根）, 则称 $f(x)$ 对于模 $p$ 是\emph{完全分裂}的. 令 $S(f)=\{\text{素数 } p\mid f(x)\ \text{对模 } p\ \text{完全分裂}\}$. 
\end{definition}

令 $\theta\in\C$ 为 $f(x)$ 的一个根, $K=\Q(\theta)$. 如果 $K/\Q$ 是伽罗瓦、Abel 或者循环扩张, 则 $f(x)$ 也分别叫作是\emph{正规}、\emph{Abel} 或者\emph{循环}多项式. 

根据上述素理想分解型式和多项式型式之间的关推论, 由定理 \ref{thm:5.2} 立即推出: 

\begin{corollary}[3]
设 $f(x),g(x)\in\Z[x]$ 均是首 $1$ 不可约正规多项式, $\alpha,\beta\in\C$ 分别是 $f(x)$ 和 $g(x)$ 的根. 如果 $S(f)$ 和 $S(g)$ 几乎相等, 则
\[
  \deg f=\deg g,
\]
并且 $\Q(\alpha)=\Q(\beta)$. 
\end{corollary}

\begin{remark}
推论 3 中关于 $f$ 和 $g$ 均为``正规''多项式这一假定是不可去掉的. Kronecker 最早研究这一问题. Gassmann 于 $1926$ 年证明存在两个 $180$ 次（！）的不可约多项式 $f$ 和 $g$, 除了有限多个素数之外对于每个素数 $p$, $f(x)$ 和 $g(x)$ 模 $p$ 均有相同的分解型式. 但是对 $f(x)$ 的每个根 $\alpha\in\C$ 和 $g(x)$ 的每个根 $\beta\in\C$, 均有 $\Q(\alpha)\neq\Q(\beta)$. $1970$ 年 I. Gerst 给出更简单的例子: $f(x)=x^8-3\cdot2^4$ 和 $g(x)=x^8-3^7$ 具有上述性质. 
\end{remark}

从推论 1 立刻推得: 

\begin{corollary}[4]
设 $f(x)\in\Z[x]$ 是首 $1$ 不可约多项式, $N$ 是 $f(x)$ 的分裂域, 则 $\delta(S(f))=1/[N:\Q]$. 
\end{corollary}

同样地可以证得: 

\begin{corollary}[5]
设 $f(x),g(x)\in\Z[x]$ 均是首 $1$ 不可约多项式. $\alpha,\beta\in\C$ 分别为 $f(x)$ 和 $g(x)$ 的根. 如果 $\Q(\alpha)=\Q(\beta)$, 则 $S(f)$ 和 $S(g)$ 几乎相等. 
\end{corollary}

\begin{proof}
事实上, 令 $K=\Q(\alpha)=\Q(\beta)$. 我们已经知道, 除了有限个之外的所有素数 $p$（即对所有 $p\nmid d(f)d(g)$ 的素数 $p$）, $p$ 在 $K$ 中素理想分解的型式与 $f(x)$ 和 $g(x)$ 模 $p$ 分解的型式均是相同的, 这就表明 $S(f)\Delta S(g)$ 是有限集合. 特别地, $S(f)$ 和 $S(g)$ 几乎相等. 
\end{proof}

\begin{example}
可以证明 $f(x)=x^4+4x^3-4x^2-40x-56$ 和 $g(x)=x^4-8x^2-24x-20$ 均是 $\Z[x]$ 中不可约多项式, 并且有相同的判别式: $d(f)=d(g)=-2^{12}\cdot3^2\cdot31$. 但是对于 $p=13$, $f(2)\equiv 0\ (\mathrm{mod}\ 13)$, 而 $g(x)$ 在 $\Z/13\Z$ 中无根. 于是由推论 5 的证明即知对于 $f(x)$ 的任一根 $\alpha\in\C$ 和 $g(x)$ 的任一根 $\beta\in\C$, 均有 $\Q(\alpha)\neq\Q(\beta)$. 
\end{example}

% ============ 第五章 §2 Abel L-函数, Чеботарёв 密度定理 (原书 pp.197-205) ============
\setcounter{example}{0}% 原书例编号分组

\section{\texorpdfstring{Abel $L$-函数, Чеботарёв 密度定理}{Abel L-函数, Chebotarev 密度定理}}\label{sec:ch5-s2}

设 $L/K$ 为数域的伽罗瓦扩张. 定理 \ref{thm:5.2} 表明, 在 $L$ 完全分裂的 $O_K$ 中素理想全体具有 $D$-密度 $1/[L:K]$. 现在我们要提出更一般的问题：对每对正整数 $f$ 和 $g$($fg=[L:K]$), 在 $L$ 上分解模式为 $(f,g)$ 的 $O_K$ 中素理想全体的 $D$-密度是多少？当 $f=1$, $g=[L:K]$ 时这就是完全分裂的情形. 

我们先对 Abel 扩张情形给出解答, 然后再考虑任意伽罗瓦扩张情形. 

\begin{theorem}\label{thm:5.3}
设 $L/K$ 是数域的 $n$ 次 Abel 扩张. 对 $G=\Gal(L/K)$ 中每个自同构 $\sigma$, 令
\[
  A(\sigma)=\left\{O_K\ \text{中素理想}\ \mathfrak{p}:\ \mathfrak{p}\ \text{在}\ L\ \text{上不分歧, 并且}\ \left(\frac{L/K}{\mathfrak{p}}\right)=\sigma\right\},
\]
则 $\delta(A(\sigma))=\frac{1}{n}$. 
\end{theorem}

\begin{proof}
映射
\[
  D_{\mathfrak{p}}\to\Gal(\overline{L}/\overline{K}),\qquad \sigma\mapsto\overline{\sigma}
\]
是群的满同态, 并且核为 $\mathfrak{p}$ 对 $L/K$ 的惯性群 $I_{\mathfrak{p}}$. 我们取 $\Gal(\overline{L}/\overline{K})$ 中自同构 $\tau\colon x\mapsto x^{|\overline{K}|}$(对 $x\in\overline{L}$)的一个原象 $\sigma_{\mathfrak{p}}\in D_{\mathfrak{p}}$, 则 $\tau$ 的全部原象为 $\sigma_{\mathfrak{p}}I_{\mathfrak{p}}$. 令 $\widehat{G}$ 为 Abel 群 $G=\Gal(L/K)$ 的特征群. 对每个特征 $\chi\in\widehat{G}$, 定义
\[
  \chi(\mathfrak{p})=\chi(\sigma_{\mathfrak{p}})\cdot\frac{1}{e(\mathfrak{P}/\mathfrak{p})}\sum_{\sigma\in I_{\mathfrak{p}}}\chi(\sigma)
  =\frac{1}{|I_{\mathfrak{p}}|}\sum_{\sigma\in I_{\mathfrak{p}}}\chi(\sigma_{\mathfrak{p}}\sigma).
\]
不难看出：

(I) $\chi(\mathfrak{p})$ 与 $\tau$ 的原象 $\sigma_{\mathfrak{p}}$ 之选取方式无关. 

(II) 若 $I_{\mathfrak{p}}$ 不包含在 $\operatorname{Ker}\chi$ 之中, 则 $\chi$ 在 $I_{\mathfrak{p}}$ 上不是平凡特征. 由特征的正交关系可知 $\chi(\mathfrak{p})=0$. 

(III) 若 $I_{\mathfrak{p}}\subseteq\operatorname{Ker}\chi$, 则 $\chi(\mathfrak{p})=\chi(\sigma_{\mathfrak{p}})$. 

(IV) 若 $\mathfrak{p}$ 在 $L$ 中不分歧, 则 $I_{\mathfrak{p}}=\{1\}$, 而 $\tau$ 只有一个原象 $\left(\frac{L/K}{\mathfrak{p}}\right)$, 从而 $\chi(\mathfrak{p})=\chi\left(\left(\frac{L/K}{\mathfrak{p}}\right)\right)$. 

现在将 $\chi$ 的定义完全积性地扩充到 $O_K$ 的全部非零整理想上, 即对于 $\mathfrak{a}=\mathfrak{p}_1\cdots\mathfrak{p}_r$, 定义 $\chi(\mathfrak{a})=\chi(\mathfrak{p}_1)\cdots\chi(\mathfrak{p}_r)$. 现在可定义一个 Dirichlet 级数
\begin{equation}\tag{1}
  L(s,\chi,L/K)=\sum_{\mathfrak{a}}\chi(\mathfrak{a})N(\mathfrak{a})^{-s},
\end{equation}
其中求和为 $\mathfrak{a}$ 过 $O_K$ 的全部非零整理想. 它叫作数域 Abel 扩张 $L/K$ 关于特征 $\chi\in\widehat{G}$ 的 \emph{Abel $L$-函数}. 这个函数有如下性质：

(a) 由于 $|\chi(\mathfrak{a})|=0$ 或 $1$, 可知级数 (1) 在 $\operatorname{Re}(s)>1$ 中定义出正则函数. 由于 $\chi$ 是完全积性的, 所以有 Euler 乘积公式
\[
  L(s,\chi,L/K)=\prod_{\mathfrak{p}}(1-\chi(\mathfrak{p})N(\mathfrak{p})^{-s})^{-1}\qquad(\operatorname{Re}(s)>1\ \text{时}).
\]

(b)
\[
  \prod_{\chi\in\widehat{G}}L(s,\chi,L/K)=\zeta_L(s)\qquad(\operatorname{Re}(s)>1),
\]
其中 $\zeta_L(s)$ 是 $L$ 的 Dedekind zeta 函数：
\[
  \zeta_L(s)=\prod_{\mathfrak{P}}(1-N(\mathfrak{P})^{-s})^{-1}\qquad(\operatorname{Re}(s)>1).
\]
由于 (b) 两边均有 Euler 展开, 所以只需对 $O_K$ 的每个素理想 $\mathfrak{p}$, 证明对应 $\mathfrak{p}$ 的因子相等：
\begin{equation}\tag{2}
  \prod_{\chi\in\widehat{G}}(1-\chi(\mathfrak{p})N(\mathfrak{p})^{-s})=\prod_{\mathfrak{B}\mid\mathfrak{p}}(1-N(\mathfrak{B})^{-s}).
\end{equation}
根据前述性质 (II) 和 (III), 可知 (2) 式左边为
\begin{equation}\tag{3}
  \prod_{\substack{\chi\in\widehat{G}\\ \operatorname{Ker}\chi\supseteq I_{\mathfrak{p}}}}(1-\chi(\mathfrak{p})N(\mathfrak{p})^{-s})
  =\prod_{\chi\in(G/I_{\mathfrak{p}})^{\wedge}}(1-\chi(\sigma_{\mathfrak{p}})N(\mathfrak{p})^{-s}).
\end{equation}
令 $\mathfrak{p}O_L=(\mathfrak{P}_1\cdots\mathfrak{P}_g)^e$, $efg=n$. 由于 $|G/I_{\mathfrak{p}}|=fg$, 而由 $\sigma_{\mathfrak{p}}$ 的取法知 $\sigma_{\mathfrak{p}}$ 在 $G/I_{\mathfrak{p}}$ 中阶为 $f$. 所以 (3) 式右边等于 $(1-N(\mathfrak{p})^{-fs})^g$. 另一方面, $\mathfrak{p}$ 的因子有 $g$ 个 $\mathfrak{P}_1,\cdots,\mathfrak{P}_g$, 而 $N(\mathfrak{P}_i)=\mathfrak{p}^f$. 于是 (2) 式右边也为 $(1-N(\mathfrak{p})^{-fs})^g$. 这就证明了 (b). 

(c) $L(s,\chi,L/K)$ 可以解析延拓到整个复平面上. 当 $\chi\neq\chi_0$ 时 $L(s,\chi,L/K)$ 在整个复平面上是正则的(没有奇点). 而 $L(s,\chi_0,L/K)=\zeta_K(s)$ 只在 $s=1$ 处有极点并且是单极点. 关于 (c) 的证明是不平凡的, 要用到类域论. 这里证明从略. 

有了以上准备, 现在可以很容易证明定理 \ref{thm:5.3}. 我们要证
\begin{equation}\tag{4}
  \lim_{s\to1^+}\sum_{\mathfrak{p}\in A(\sigma)}N(\mathfrak{p})^{-s}\big/-\log(s-1)=\frac{1}{n}.
\end{equation}
令
\[
  T(s)=\frac{1}{n}\sum_{\chi\in\widehat{G}}\chi(\sigma^{-1})\log L(s,\chi,L/K).
\]
由上面的 (b) 和 (c), 可知当 $\chi\neq\chi_0$ 时, $L(1,\chi,L/K)\neq 0,\infty$. 因此, 
\[
  \lim_{s\to1^+}T(s)/-\log(s-1)
  =\lim_{s\to1^+}\frac{1}{n}\cdot\frac{\log\zeta_K(s)}{-\log(s-1)}=\frac{1}{n}.
\]
但是
\begin{align*}
  T(s)&=\frac{1}{n}\sum_{\chi\in\widehat{G}}\chi(\sigma^{-1})\log\prod_{\mathfrak{p}}(1-\chi(\mathfrak{p})N(\mathfrak{p})^{-s})^{-1}\\
  &=\frac{1}{n}\sum_{\chi\in\widehat{G}}\chi(\sigma^{-1})\sum_{\mathfrak{p}}\log(1-\chi(\mathfrak{p})N(\mathfrak{p})^{-s})^{-1}\\
  &=\frac{1}{n}\sum_{\chi\in\widehat{G}}\chi(\sigma^{-1})\sum_{\mathfrak{p}}\sum_{m=1}^{\infty}\frac{\chi(\mathfrak{p})^mN(\mathfrak{p})^{-ms}}{m}\\
  &=\frac{1}{n}\sum_{\chi\in\widehat{G}}\chi(\sigma^{-1})\sum_{\mathfrak{p}}\chi(\mathfrak{p})N(\mathfrak{p})^{-s}+g(s)\\
  &=\frac{1}{n}\sum_{\mathfrak{p}}N(\mathfrak{p})^{-s}\sum_{\chi\in\widehat{G}}\chi(\sigma^{-1})\chi(\mathfrak{p})+g(s),
\end{align*}
其中 $g(s)$ 是对 $m\geqslant 2$ 求和部分, 如同以前一样可知 $g(s)$ 在 $s=1$ 附近是有界的. 而除了有限个 $\mathfrak{p}$ 之外, 均有 $\chi(\mathfrak{p})=\chi\left(\left(\frac{L/K}{\mathfrak{p}}\right)\right)$. 由特征正交关系知
\[
  \sum_{\chi\in\widehat{G}}\chi(\sigma^{-1})\chi\left(\left(\frac{L/K}{\mathfrak{p}}\right)\right)
  =\begin{cases}
    n, & \text{若}\ \sigma=\left(\frac{L/K}{\mathfrak{p}}\right),\\
    0, & \text{否则}.
  \end{cases}
\]
这就表明
\[
  T(s)=\sum_{\mathfrak{p}\in A(\sigma)}N(\mathfrak{p})^{-s}+g'(s),
\]
其中 $g'(s)$ 在 $s=1$ 附近有界. 综合上述可知
\[
  \frac{1}{n}=\lim_{s\to1^+}T(s)/-\log(s-1)
  =\lim_{s\to1^+}\sum_{\mathfrak{p}\in A(\sigma)}N(\mathfrak{p})^{-s}/-\log(s-1)
  =\delta(A(\sigma)).
\]
这就证明了定理 \ref{thm:5.3}. 
\end{proof}

\setcounter{corollary}{0}% 原书第五章系按节重新编号

\begin{corollary}[1]
设 $L/K$ 是 Abel 扩张, $fg=n=[L:K]$. 以 $n_f$ 表示 $G=\Gal(L/K)$ 中 $f$ 阶元素的个数, 而
\[
  A=\left\{O_K\ \text{中素理想}\ \mathfrak{p}:\ \mathfrak{p}\ \text{在}\ L\ \text{中素理想分解型式为}\ (f,g)\right\},
\]
则 $\delta(A)=\frac{n_f}{n}$. 
\end{corollary}

\begin{example}
设 $L/K$ 是数域的 Abel 扩张, 并且 $G=\Gal(L/K)$ 为两个 $2$ 阶循环群的直积, 则 $G$ 中有一个 $1$ 阶元素和 $3$ 个 $2$ 阶元素. 所以 $O_K$ 中满足 $\mathfrak{p}O_L=\mathfrak{P}_1\mathfrak{P}_2$ 的全部素理想 $\mathfrak{p}$ 其 $D$-密度为 $\frac{3}{4}$. 而在 $L$ 中完全分裂的全部 $\mathfrak{p}$ 具有 $D$-密度 $\frac{1}{4}$. 最后, 在 $L$ 中惯性的(即 $f=4$)全部 $\mathfrak{p}$ 具有 $D$-密度 $0$, 因为 $G$ 中没有 $4$ 阶元素. 注意：一个 $D$-密度为 $0$ 的集合不必为空集, 它可以是有限集合甚至是无限集合. 但事实上对于本例中的情形, 在 $L$ 中惯性的 $\mathfrak{p}$ 是不存在的. (为什么？)
\end{example}

\begin{remark}
(1) 由定理 \ref{thm:5.3} 的 (c) 可知
\[
  \zeta_L(s)/\zeta_K(s)=\prod_{\chi\neq\chi_0}L(s,\chi,L/K)
\]
是整个复平面上的正则函数. 换句话说, (利用类域论)可以证明 Artin 猜想对于 $L/K$ 为 Abel 扩张的情形是成立的. 

(2) 定理 \ref{thm:5.1} 是定理 \ref{thm:5.3} 的一个特例. 因为在定理 \ref{thm:5.3} 中取 $K=\Q$, $L=\Q(\zeta_k)$, 则 $[L:\Q]=\varphi(k)$, 而
\[
  G=\Gal(L/\Q)=\{\sigma_l:(l,k)=1,\ 1\leqslant l<k\},
\]
其中 $\sigma_l(\zeta_k)=\zeta_k^l$. 若 $p$ 在 $L$ 中不分歧(即 $p\nmid k$), 我们已经知道 $\left(\frac{L/K}{p}\right)=\sigma_p$, 所以 $\left(\frac{L/K}{p}\right)=\sigma_l\Leftrightarrow p\equiv l\ (\mathrm{mod}\ k)$. 所以素数集合 $A=\{p:p\equiv l\ (\mathrm{mod}\ k)\}$ 的 $D$-密度就是素数集合 $A'=\left\{p:\left(\frac{L/K}{p}\right)=\sigma_l\right\}$ 的 $D$-密度. 根据定理 \ref{thm:5.3}, $\delta(A')=\frac{1}{\varphi(k)}$, 从而 $\delta(A)=\frac{1}{\varphi(k)}$. 

(3) 我们略去了定理 \ref{thm:5.2} 的一部分证明, 因为证明 Abel $L$-函数解析延拓和奇点特性时要用到类域论. 但是当 $K=\Q$ 时, 我们可以不用类域论而只用 Kronecker-Weber 定理证明定理 \ref{thm:5.3}. 因为若 $L/\Q$ 是 Abel 扩张, 则 $L$ 是分圆域 $\Q(\zeta_m)$ 的子域, 其中 $m$ 为域 $L$ 的导子. 伽罗瓦群 $\Gal(\Q(\zeta_m)/\Q)$ 自然同构于 $(\Z/m\Z)^{\times}$. 设 $N$ 为 $L$ 的固定子群, $N$ 在 $(\Z/m\Z)^{\times}$ 中的象为 $\varphi(N)$, 则 $G=\Gal(L/\Q)$ 自然同构于 $(\Z/m\Z)^{\times}/\varphi(N)$(参见定理 \ref{thm:2.18} 的证明和后面的注记 1). 若素数 $p$ 在 $L$ 中不分歧, 我们不妨设 $(p,m)=1$(因为 $m$ 的有限个素因子是 $D$-密度为 $0$ 的集合, 不影响以后的讨论). $G$ 中每个元素 $\sigma$(即 $L$ 的自同构)自然看成是 $(\Z/m\Z)^{\times}$ 模 $\varphi(N)$ 的一个陪集, 所以相当于模 $m$ 的 $|\varphi(N)|$ 个同余类. 注意 $|\varphi(N)|=|N|=[\Q(\zeta_m):L]=\frac{\varphi(m)}{[L:\Q]}$. 如果 $\left(\frac{L/\Q}{p}\right)=\sigma_{\mathfrak{p}}$, 所以这相当于 $p$ 属于 $\sigma$ 所对应的那 $\frac{\varphi(m)}{[L:\Q]}$ 个模 $m$ 同余类. 根据例 1, 每个同余类中素数的 $D$-密度为 $\frac{1}{\varphi(m)}$. 所以满足 $\left(\frac{L/\Q}{p}\right)=\sigma$ 的所有素数组成的集合, 其 $D$-密度为 $\frac{\varphi(m)}{[L:\Q]}\cdot\varphi(m)^{-1}=[L:\Q]^{-1}$. 这就是定理 \ref{thm:5.3}(对于 $K=\Q$ 的情形). 
\end{remark}

现在讨论 $L/K$ 为伽罗瓦扩张的情形. 这时, 设 $O_K$ 中素理想 $\mathfrak{p}$ 在 $L$ 中不分歧, 则 $\mathfrak{p}O_L=\mathfrak{P}_1\cdots\mathfrak{P}_g$, 令 $\mathfrak{P}=\mathfrak{P}_1$, 则每个 $\mathfrak{P}_i$ 均可写成 $\sigma(\mathfrak{P})$(对某个 $\sigma\in G=\Gal(L/K)$). 但是
\[
  \left(\frac{L/K}{\sigma(\mathfrak{P})}\right)=\sigma\left(\frac{L/K}{\mathfrak{P}}\right)\sigma^{-1},
\]
于是 $\left\{\left(\frac{L/K}{\mathfrak{P}_i}\right)\mid 1\leqslant i\leqslant g\right\}=\left\{\sigma\left(\frac{L/K}{\mathfrak{P}}\right)\sigma^{-1}\mid\sigma\in G\right\}$, 这是有限群 $G$ 的一个共轭类, 我们将这个共轭类表示成 $\left(\frac{L/K}{\mathfrak{p}}\right)$(当 $L/K$ 为 Abel 扩张时, 这个共轭类只含有一个元素, 我们曾经把这个元素表示成 $\left(\frac{L/K}{\mathfrak{p}}\right)$). 

\begin{theorem}[\textbf{Чеботарёв}]\label{thm:5.4}
设 $L/K$ 是数域的伽罗瓦扩张, $G=\Gal(L/K)$. $C$ 为 $G$ 中一个共轭元素类, $|C|=c$. $A=\left\{O_K\ \text{中素理想}\ \mathfrak{p}\mid\mathfrak{p}\ \text{在}\ L\ \text{中不分歧并且}\ \left(\frac{L/K}{\mathfrak{p}}\right)=C\right\}$, 则 $\delta(A)=c/n$, $n=[L:K]$. 
\end{theorem}

\begin{proof}
俄国数学家 Чеботарёв 原来的证明很复杂. 这里采用 $1969$ 年 McCluer 给出的简化证明. 首先需要下列引理. 

\setcounter{lemma}{4}% 原书第五章引理编号跳号(无引理2-4)

\begin{lemma}\label{lem:ch5-5}
设 $L/K$ 是数域的伽罗瓦扩张, $\mathfrak{P}$ 和 $\mathfrak{p}$ 分别是 $O_L$ 和 $O_K$ 中的(固定)素理想, $\mathfrak{P}\mid\mathfrak{p}$, $\mathfrak{p}$ 在 $L$ 中不分歧. $M$ 是 $L/K$ 的中间域. 并且对于 $O_M$ 中每个素理想 $\mathfrak{q}\mid\mathfrak{p}$, $\mathfrak{q}$ 在 $L$ 中均是惯性的(从而 $\mathfrak{q}O_L$ 为 $O_L$ 中素理想). 令 $\left(\frac{L/K}{\mathfrak{P}}\right)=\sigma\in G=\Gal(L/K)$, 则 $\mathfrak{p}$ 在 $O_M$ 中共存在 $[C_G(\sigma):H]$ 个素理想因子 $\mathfrak{q}$, 使得 $\left(\frac{L/K}{\mathfrak{q}O_L}\right)=\sigma$. 其中 $C_G(\sigma)=\{\tau\in G\mid\tau\sigma=\sigma\tau\}$($\sigma$ 在 $G$ 中的中心化子), 而 $H=\langle\sigma\rangle$($\sigma$ 在 $G$ 中生成的循环子群). 
\end{lemma}

这是因为：从引理的假设可知, 我们只需证明 $O_L$ 中共有 $[C_G(\sigma):H]$ 个 $\mathfrak{P}'\mid\mathfrak{p}$ 使得 $\left(\frac{L/K}{\mathfrak{P}'}\right)=\sigma$. 每个这种 $\mathfrak{P}'$ 均可写成 $\tau(\mathfrak{P})$, $\tau\in G$. 而
\[
  \left(\frac{L/K}{\tau(\mathfrak{P})}\right)=\sigma
  \Leftrightarrow\tau\sigma\tau^{-1}=\sigma
  \Leftrightarrow\tau\in C_G(\sigma);
\]
并且
\[
  \tau_1\mathfrak{P}=\tau_2\mathfrak{P}
  \Leftrightarrow\tau_2^{-1}\tau_1\mathfrak{P}=\mathfrak{P}
  \Leftrightarrow\tau_2^{-1}\tau_1\in D_{\mathfrak{P}}=\langle\sigma\rangle=H.
\]
由此即得引理. 

现在证明定理 \ref{thm:5.4}：这只需证明
\[
  \sum_{\left(\frac{L/K}{\mathfrak{p}}\right)=C}N(\mathfrak{p})^{-s}
  =-\frac{c}{n}\log(s-1)+o(\log(s-1))\qquad(\text{当 } s\to1^+\ \text{时}).
\]
取 $\sigma\in C$, $H=\langle\sigma\rangle\subseteq G=\Gal(L/K)$. $M$ 为 $H$ 的固定子域, 则 $\Gal(L/M)\cong H$, 从而 $L/M$ 为循环扩张. 由定理 \ref{thm:5.3} 知
\[
  \sum_{\left(\frac{L/M}{\mathfrak{q}}\right)=\sigma}N(\mathfrak{q})^{-s}
  =-\frac{1}{|H|}\log(s-1)+o(\log(s-1)).
\]
从而
\[
  \sum_{\substack{\mathfrak{q}\mid\mathfrak{p}\\ \left(\frac{L/M}{\mathfrak{q}}\right)=\sigma\\ f(\mathfrak{q}\mid\mathfrak{p})=1}}N(\mathfrak{p})^{-s}
  =-\frac{1}{|H|}\log(s-1)+o(\log(s-1)).
\]
对于使 $\left(\frac{L/M}{\mathfrak{q}}\right)=\sigma$ 和 $\mathfrak{q}\mid\mathfrak{p}$ 成立的每个 $\mathfrak{p}$, $M$ 满足上面引理条件. 因为 $\left(\frac{L/M}{\mathfrak{q}}\right)=\sigma$, 则分解群 $D_{\mathfrak{p}}=\langle\sigma\rangle\cong\Gal(L/M)$, 从而 $\mathfrak{q}$ 在 $L$ 中惯性. 若 $\mathfrak{q}'$ 为 $M$ 中素理想并且 $\mathfrak{q}'\mid\mathfrak{p}$, 则 $D_{\mathfrak{q}'}$ 与 $D_{\mathfrak{q}}$ 共轭, 从而仍有 $|D_{\mathfrak{q}'}|=|\Gal(L/M)|$, 则 $\mathfrak{q}'$ 在 $L$ 中惯性. 于是由引理得出(注意 $\left(\frac{L/M}{\mathfrak{q}}\right)=\sigma\Leftrightarrow\left(\frac{L/K}{\mathfrak{p}}\right)=C$)
\[
  [C_G(\sigma):H]\sum_{\left(\frac{L/K}{\mathfrak{p}}\right)=C}N(\mathfrak{p})^{-s}
  =-\frac{1}{|H|}\log(s-1)+o(\log(s-1)).
\]
于是
\[
  \sum_{\left(\frac{L/K}{\mathfrak{p}}\right)=C}N(\mathfrak{p})^{-s}
  =-\frac{1}{[C_G(\sigma):H]\cdot|H|}\log(s-1)+o(\log(s-1)).
\]
但是 $[C_G(\sigma):H]\cdot|H|=|C_G(\sigma)|=\frac{n}{c}$, 这就表明 $\delta(A)=\frac{c}{n}$. 
\end{proof}

我们可以把定理 \ref{thm:5.3} 的系 1 推广到任意伽罗瓦扩张的情形. 

\begin{corollarynonum}
设 $L/K$ 是数域的伽罗瓦扩张, $fg=n=[L:K]$. 以 $n_f$ 表示 $G=\Gal(L/K)$ 中 $f$ 阶元素的个数, 令
\[
  A=\{O_K\ \text{的素理想}\ \mathfrak{p}:\mathfrak{p}\ \text{在}\ L\ \text{中分解型式为}\ (f,g)\},
\]
则 $\delta(A)=\frac{n_f}{n}$. 
\end{corollarynonum}

\begin{proof}
设 $\mathfrak{p}$ 在 $L$ 中不分歧, 则 $\mathfrak{p}$ 在 $L$ 中分解型式为 $(f,g)$ 当且仅当 $G$ 的共轭元素类 $\left(\frac{L/K}{\mathfrak{p}}\right)$ 中所有元素的阶均为 $f$. 再由定理 \ref{thm:5.4} 即得此系. 
\end{proof}

% ============ 习题 ============
\begin{Exercises}
  \item 设 $f(x)$ 和 $g(x)$ 为 $\Z[x]$ 中不可约多项式, $\alpha$ 和 $\beta$ 分别是 $f(x)$ 和 $g(x)$ 在复数域中的一个根. 如果 $\Q(\alpha)\subseteq\Q(\beta)$, 求证 $S(f)$ 几乎包含 $S(g)$, 即素数集合 $\{p:p\in S(g),\ p\notin S(f)\}$ 的 $D$-密度为 $0$. 
  \item 设 $f_1(x),\cdots,f_r(x)$ 是 $\Z[x]$ 中有限个多项式, 并且次数均 $\geqslant 1$. 求证存在无限多个素数 $p$, 使得 $f_1(x),\cdots,f_r(x)$ 模 $p$ 均完全分裂. 
  \item 设 $L/K$ 是数域的伽罗瓦扩张. 求证：$O_K$ 中有无限多个素理想 $\mathfrak{p}$ 在 $L$ 中惯性的充分必要条件为 $L/K$ 是循环扩张(即伽罗瓦群 $\Gal(L/K)$ 为循环群). 
  \item 计算素数集合 $A=\{p:2$ 为模 $p$ 的三次剩余$\}$ 和 $A'=\{p:2$ 为模 $p$ 的四次剩余$\}$ 的 $D$-密度. 
  \item 设 $f(x)$ 是 $\Z[x]$ 中次数 $\geqslant 2$ 的首 $1$ 不可约多项式, 求证存在无限多素数 $p$, 使得 $f(x)$ 在有限域 $\F_p$ 中无根. 
\end{Exercises}

